\documentclass[11pt]{article}

\usepackage[margin=1in]{geometry}
\usepackage[T1]{fontenc}
\usepackage{lmodern}
\usepackage{amsmath}
\usepackage{booktabs}
\usepackage{enumitem}
\usepackage{hyperref}
\usepackage{xcolor}
\usepackage{listings}

\hypersetup{
    colorlinks=true,
    linkcolor=blue,
    urlcolor=blue,
    citecolor=blue
}

\lstset{
    basicstyle=\ttfamily\small,
    breaklines=true,
    frame=single,
    columns=fullflexible,
    showstringspaces=false
}

\title{ADR-003\\Constrained Policy Template Parameters}
\author{William Gordon Carter}
\date{\today}

\begin{document}

\maketitle

\begin{center}
\begin{tabular}{ll}
\toprule
\textbf{Field} & \textbf{Value} \\
\midrule
ADR & 003 \\
Title & Constrained Policy Template Parameters \\
Status & Proposed \\
Author & William Gordon Carter \\
Created & \today \\
Supersedes & None \\
Superseded By & --- \\
Depends On & ADR-001 Compile-Time Policy Architecture \\
Related & ADR-002 Compile-Time Environment Policy Architecture \\
\bottomrule
\end{tabular}
\end{center}

\tableofcontents
\newpage

\section{Context}

ADR-001 establishes NavKit's broad compile-time policy architecture:

\begin{center}
Concept $\rightarrow$ CRTP Base $\rightarrow$ Concrete Policy $\rightarrow$ Runtime Algorithm
\end{center}

As this architecture is extended into the estimator stack, class templates such as
\texttt{KalmanFilter}, \texttt{Sensor}, and \texttt{Navigator} need to express
their compile-time policy requirements clearly.

The desired API should make template parameters read like architectural roles
rather than unconstrained implementation types.

\section{Decision}

NavKit shall prefer constrained policy template parameters in public class and
function template declarations.

For unary concepts:

\begin{lstlisting}[language=C++]
template<planet::PlanetPolicy Planet>
struct Spherical;
\end{lstlisting}

For concepts that depend on a previously declared template parameter, NavKit
shall define the candidate type as the first concept parameter. This enables
the clean constrained-parameter syntax:

\begin{lstlisting}[language=C++]
template<
    StateDefPolicy StateDef,
    InjectionPolicy<StateDef> Injection,
    ResetPolicy<StateDef> Reset>
class KalmanFilter;
\end{lstlisting}

This style is preferred over exposing unconstrained \texttt{typename}
parameters with a separate \texttt{requires} clause when the constrained
parameter syntax is available and readable.

\section{C++ Concept Parameter Ordering}

For constrained template parameter syntax to work cleanly, the candidate type
must be the first template parameter of the concept.

For example:

\begin{lstlisting}[language=C++]
template<typename T, typename StateDef>
concept InjectionPolicy =
    StateDefPolicy<StateDef> &&
    requires(State<StateDef>& x, const State<StateDef>& dx)
{
    { T::apply(x, dx) } -> std::same_as<void>;
};
\end{lstlisting}

The concept definition itself uses unconstrained \texttt{typename} parameters
and composes any prerequisite concepts in the concept body. In standard C++,
concept definitions cannot have associated constraints on their own template
parameter list. For example, this tempting spelling is invalid:

\begin{lstlisting}[language=C++]
template<typename T, StateDefPolicy StateDef>
concept InjectionPolicy = ...; // invalid concept definition syntax
\end{lstlisting}

The constrained-parameter syntax is reserved for templates that consume the
concept, not for the definition of the concept itself.

Then this declaration:

\begin{lstlisting}[language=C++]
template<
    StateDefPolicy StateDef,
    InjectionPolicy<StateDef> Injection>
class Example;
\end{lstlisting}

is interpreted as:

\begin{lstlisting}[language=C++]
template<
    typename StateDef,
    typename Injection>
requires StateDefPolicy<StateDef> &&
         InjectionPolicy<Injection, StateDef>
class Example;
\end{lstlisting}

The constrained parameter \texttt{Injection} is implicitly inserted as the
first argument to the concept.

Therefore, NavKit concepts intended for constrained template parameter syntax
shall follow this ordering:

\begin{lstlisting}[language=C++]
template<typename Candidate, typename Context>
concept SomePolicy = ...;
\end{lstlisting}

not:

\begin{lstlisting}[language=C++]
template<typename Context, typename Candidate>
concept SomePolicy = ...;
\end{lstlisting}

and not:

\begin{lstlisting}[language=C++]
template<typename Candidate, SomeOtherPolicy Context>
concept SomePolicy = ...;
\end{lstlisting}

\section{Design Rationale}

\subsection{Readable Architecture}

The preferred form:

\begin{lstlisting}[language=C++]
template<
    StateDefPolicy StateDef,
    InjectionPolicy<StateDef> Injection,
    ResetPolicy<StateDef> Reset>
class KalmanFilter;
\end{lstlisting}

reads like an architectural declaration:

\begin{quote}
A Kalman filter is parameterized by a state definition, an injection policy
compatible with that state definition, and a reset policy compatible with that
state definition.
\end{quote}

This is more expressive than:

\begin{lstlisting}[language=C++]
template<
    typename StateDef,
    typename Injection,
    typename Reset>
requires StateDefPolicy<StateDef> &&
         InjectionPolicy<Injection, StateDef> &&
         ResetPolicy<Reset, StateDef>
class KalmanFilter;
\end{lstlisting}

Both are valid. The constrained-parameter form is preferred when it remains
clear.

\subsection{Concepts as Policy Interfaces}

Concepts define the compile-time policy interface. A concrete policy does not
explicitly declare that it implements a concept. Instead, the compiler checks
whether the concrete policy satisfies the required expressions.

For example:

\begin{lstlisting}[language=C++]
struct AdditiveInjection
{
    static void apply(State<MyStateDef>& x, const State<MyStateDef>& dx);
};
\end{lstlisting}

satisfies \texttt{InjectionPolicy<AdditiveInjection, MyStateDef>} if the
required \texttt{apply} expression is valid.

\subsection{Compatibility is Part of the Policy Concept}

For state-dependent policies, compatibility with the selected state definition
is part of the policy concept itself.

That is, \texttt{InjectionPolicy<Injection, StateDef>} means both:

\begin{enumerate}
    \item \texttt{Injection} has the required injection-policy interface.
    \item \texttt{Injection} can operate on \texttt{State<StateDef>}.
\end{enumerate}

A separate \texttt{CompatibleInjectionPolicy} concept may be introduced later
if needed, but it is not required for the initial estimator refactor.

\section{Kalman Filter Target Architecture}

The target Kalman filter declaration is:

\begin{lstlisting}[language=C++]
template<
    StateDefPolicy StateDef,
    InjectionPolicy<StateDef> Injection,
    ResetPolicy<StateDef> Reset>
class KalmanFilter;
\end{lstlisting}

The filter remains a runtime algorithm. It is not itself a policy.

Its template parameters are compile-time policies that configure the runtime
algorithm.

\section{Injection Policy}

An injection policy defines how the accumulated error state is injected into
the nominal state.

Example concept:

\begin{lstlisting}[language=C++]
template<typename T, typename StateDef>
concept InjectionPolicy =
    StateDefPolicy<StateDef> &&
    requires(State<StateDef>& x, const State<StateDef>& dx)
{
    { T::apply(x, dx) } -> std::same_as<void>;
};
\end{lstlisting}

Example concrete policy:

\begin{lstlisting}[language=C++]
struct AdditiveInjection
{
    static void apply(State<MyStateDef>& x, const State<MyStateDef>& dx)
    {
        x.vector() -= dx.vector();
    }
};
\end{lstlisting}

If an injection policy is only valid for a particular configured state
definition, that constraint should be expressed by the function signature or by
requires-clauses inside the policy.

\section{Reset Policy}

A reset policy defines how the error state and covariance are reset after
injection.

Example concept:

\begin{lstlisting}[language=C++]
template<typename T, typename StateDef>
concept ResetPolicy =
    StateDefPolicy<StateDef> &&
    requires(State<StateDef>& x,
             State<StateDef>& dx,
             StateCov<StateDef>& P)
{
    { T::reset_dx(dx) } -> std::same_as<void>;
    { T::reset_covariance(x, dx, P) } -> std::same_as<void>;
};
\end{lstlisting}

Example target use:

\begin{lstlisting}[language=C++]
using Filter = KalmanFilter<
    NavStateDef,
    AdditiveInjection,
    DefaultReset>;
\end{lstlisting}

\section{Measurement Policy}

Measurement policies should follow the same ordering rule.

Example:

\begin{lstlisting}[language=C++]
template<typename T, typename StateDef>
concept MeasurementPolicy =
    StateDefPolicy<StateDef> &&
    requires(const State<StateDef>& x, const typename T::NoiseContext& ctx)
{
    T::M;
    typename T::O_t;
    typename T::R_t;
    typename T::H_t;
    typename T::K_t;
    typename T::NoiseContext;

    { T::obs(x) } -> std::same_as<typename T::O_t>;
    { T::compute_h(x) } -> std::same_as<typename T::H_t>;
    { T::compute_r(ctx) } -> std::same_as<typename T::R_t>;
} &&
requires
{
    requires T::M > 0;
    requires T::O_t::RowsAtCompileTime == T::M;
    requires T::O_t::ColsAtCompileTime == 1;
    requires T::R_t::RowsAtCompileTime == T::M;
    requires T::R_t::ColsAtCompileTime == T::M;
    requires T::H_t::RowsAtCompileTime == T::M;
    requires T::H_t::ColsAtCompileTime == StateDef::N;
    requires T::K_t::RowsAtCompileTime == StateDef::N;
    requires T::K_t::ColsAtCompileTime == T::M;
};
\end{lstlisting}

\texttt{K\_t} is included in the measurement-model boundary because the current
\texttt{KalmanFilter} records and exposes Kalman-gain statistics alongside
innovation, covariance, and Jacobian diagnostics.

This enables declarations such as:

\begin{lstlisting}[language=C++]
template<
    StateDefPolicy StateDef,
    MeasurementPolicy<StateDef> Measurement>
void observation_update(...);
\end{lstlisting}

\section{Sensor and Noise Policies}

Sensor noise policies operate at a narrower boundary than full measurement
models. A \texttt{Sensor} needs a model-provided noise context and a concrete
measurement sample type; it does not know the state definition. Full
state/model compatibility remains checked where a filter consumes a sensor.

For example:

\begin{lstlisting}[language=C++]
template<typename T, typename Model, typename MeasurementSample>
concept NoisePolicy =
    requires(typename Model::NoiseContext& ctx,
             const MeasurementSample& meas)
{
    { T::update(ctx, meas) } -> std::same_as<void>;
};
\end{lstlisting}

This enables declarations such as:

\begin{lstlisting}[language=C++]
template<
    typename Model,
    std::size_t BufferSize,
    NoisePolicy<Model, Measurement<Model::M>> Noise>
class Sensor;
\end{lstlisting}

Where the constrained-parameter syntax becomes unclear, NavKit may use a
standard \texttt{requires} clause. Readability is the guiding criterion.

\section{Navigator Target Architecture}

The Navigator should remain an orchestration runtime algorithm.

A future target declaration is:

\begin{lstlisting}[language=C++]
template<
    FilterPolicy Filter,
    SensorCollectionPolicy Sensors,
    PropagationPolicy Propagation,
    UpdatePolicy Update>
class Navigator;
\end{lstlisting}

The Navigator is not directly parameterized by planet, gravity, or frame
policies. Those are expected to be encoded inside the selected propagation or
mechanization policy.

For example:

\begin{lstlisting}[language=C++]
using Planet = planet::Wgs84;
using Gravity = gravity::J2<Planet>;
using Propagation = mechanization::Pcpf<Planet, Gravity>;

using Nav = Navigator<
    Filter,
    Sensors,
    Propagation,
    UpdatePolicy>;
\end{lstlisting}

The concrete \texttt{Navigator} type indirectly encodes the environment, but
the Navigator interface only depends on the propagation policy concept.

\section{Static Assertions and Compile-Time Tests}

Every compile-time policy architecture should have compile-time unit tests that
demonstrate both valid and invalid configurations.

Example:

\begin{lstlisting}[language=C++]
static_assert(StateDefPolicy<NavStateDef>);
static_assert(InjectionPolicy<AdditiveInjection, NavStateDef>);
static_assert(ResetPolicy<DefaultReset, NavStateDef>);

static_assert(!InjectionPolicy<BadInjection, NavStateDef>);
\end{lstlisting}

Negative compile-time tests should generally use \texttt{static\_assert(!...)}
rather than intentionally creating code that fails to compile.

This allows the test suite to document invalid configurations while still
compiling successfully.

\section{Boundary Checks}

Concepts and \texttt{static\_assert}s should be used at policy boundaries,
where user-selected or framework-selected template types first enter an
algorithm.

Examples:

\begin{itemize}
    \item \texttt{KalmanFilter} checks \texttt{StateDefPolicy},
    \texttt{InjectionPolicy}, and \texttt{ResetPolicy}.
    \item \texttt{MechanizationPcpf} checks \texttt{PlanetPolicy},
    \texttt{GravityPolicy}, and propagation requirements.
    \item \texttt{Sensor} checks measurement and noise policies.
\end{itemize}

The goal is to produce clear diagnostics at the architectural boundary rather
than allowing invalid types to fail deep inside matrix operations or nested
template instantiations.

\section{Naming Convention}

NavKit shall use the following naming convention:

\begin{center}
\begin{tabular}{ll}
\toprule
\textbf{Construct} & \textbf{Name Pattern} \\
\midrule
Concept & \texttt{<Domain>Policy} \\
CRTP base & \texttt{<Domain>PolicyBase} \\
Concrete policy & Descriptive implementation name \\
Runtime algorithm & Descriptive algorithm name \\
\bottomrule
\end{tabular}
\end{center}

Examples:

\begin{itemize}
    \item \texttt{PlanetPolicy}
    \item \texttt{PlanetPolicyBase}
    \item \texttt{Wgs84}
    \item \texttt{GravityPolicy}
    \item \texttt{GravityPolicyBase}
    \item \texttt{J2}
    \item \texttt{InjectionPolicy}
    \item \texttt{AdditiveInjection}
\end{itemize}

The term \texttt{Concept} shall not be used in type names. The concept
represents the policy interface, so the domain role is the important name.

\section{Alternatives Considered}

\subsection{Unconstrained Template Parameters}

Using \texttt{typename} for every template parameter was rejected as the
primary style because it hides policy requirements from the declaration and
produces poorer diagnostics.

\subsection{Requires Clauses Everywhere}

Explicit \texttt{requires} clauses are valid and sometimes necessary. However,
when constrained-parameter syntax is readable, it better communicates the
architecture.

\subsection{Separate Compatibility Concepts}

Separate compatibility concepts may be useful for complex relationships, but
they are not required for the initial Kalman filter and Navigator refactor.
For now, state compatibility is part of the binary policy concept.

\section{Consequences}

\subsection{Benefits}

\begin{itemize}
    \item Template declarations read like architecture.
    \item Compiler diagnostics improve.
    \item Policy requirements are visible at API boundaries.
    \item Invalid configurations can be tested with \texttt{static\_assert}.
    \item The estimator stack aligns with ADR-001.
\end{itemize}

\subsection{Costs}

\begin{itemize}
    \item Concepts must be carefully ordered with candidate type first.
    \item Some complex constraints may still require explicit \texttt{requires} clauses.
    \item Refactoring existing unconstrained templates requires care.
\end{itemize}

\section{Decision Summary}

NavKit shall use constrained policy template parameters wherever practical.
Concept definitions shall place the candidate policy type first so dependent
concept syntax such as \texttt{InjectionPolicy<StateDef> Injection} can be used
in class and function template declarations.

This ADR establishes the template declaration style for the Kalman filter,
sensor, measurement, reset, injection, propagation, and Navigator architecture
refactors that precede full INS mechanization.

\end{document}
